\documentclass[10pt,letterpaper]{article}
\usepackage{cornell-notes}

\title{Clause 01: Scope}
\author{}
\date{}

\setDocTitle{Clause 01: Scope}
\setDocSubtitle{{C 2024}}
\setDocOwner{{C 2024 Cornell Notes}}
\setDocDate{{\today}}
\setCornellCollection{{C 2024 Cornell Notes}}
\setCornellUnitType{{Clause}}
\setCornellUnitNumber{01}
\setCornellUnitTitle{Scope}
\hypersetup{pdftitle={Clause 01: Scope},pdfsubject={C 2024 study notes},pdfauthor={}}

\begin{document}
\maketitle



\begin{CornellOverviewBox}{Clause Orientation}
This clause establishes what the source material specifies and what it deliberately leaves to implementations, host systems, and external processes. It frames the specification as a portability contract for both programmers and implementers: program representation, language rules, data representation, and implementation limits are in scope, while concrete translation, invocation, and external data-transfer mechanisms are not.
\end{CornellOverviewBox}

\section{Learning Objectives}
\begin{itemize}
  \item Identify the six categories of subject matter governed by the source material.
  \item Distinguish specified program behavior from mechanisms intentionally left outside the specification.
  \item Explain how the scope supports portability across different data-processing systems.
  \item Separate language and implementation limits from the capacity of a particular machine.
  \item Use the scope as a boundary check when evaluating a portability claim.
\end{itemize}

\section{Normative Structure Map}
\begin{CornellNotesTable}
\CornellNoteRow{Scope boundary}{The clause defines the line between the portable language contract and the implementation environment.}
\CornellNoteRow{Explicit exclusions}{Translation mechanisms, machine capacity, and external process details are outside the language specification.}
\CornellNoteRow{Portability contract}{The scope guides questions about language meaning, representation, and required implementation limits.}
\end{CornellNotesTable}

\section{What the Specification Governs}
\begin{CornellNotesTable}
\CornellNoteRow{What is the central contract?}{The source defines the form and interpretation of programs written in the language. Its purpose is portable meaning: conforming source should be analyzable against one set of syntactic, semantic, and environmental rules.}
\CornellNoteRow{Which program properties are specified?}{\begin{itemize}[leftmargin=*,nosep]
  \item Representation of programs
  \item Syntax and constraints
  \item Semantic interpretation
  \item Representation of program input and output
  \item Restrictions and minimum limits imposed by conforming implementations
\end{itemize}}
\CornellNoteRow{Who uses this boundary?}{Programmers use it to reason about portable source; implementers use it to determine what translators, execution environments, and libraries must accept or document.}
\CornellNoteRow{Why is portability the organizing goal?}{The rules define a shared abstract contract across systems whose processors, operating environments, encodings, and resource capacities may differ.}
\end{CornellNotesTable}

\begin{CornellSummaryBox}{Section Summary}
The clause governs the meaning and portable representation of programs, not the engineering details of any one compiler, linker, loader, operating system, or processor.
\end{CornellSummaryBox}



\section{What the Specification Does Not Govern}
\begin{CornellNotesTable}
\CornellNoteRow{Is a translation mechanism prescribed?}{No. The required behavior of translation is constrained, but the concrete mechanism used to transform a program is outside scope.}
\CornellNoteRow{Is program invocation prescribed?}{No. How an executable is launched is a host or implementation matter unless a later requirement constrains observable program behavior.}
\CornellNoteRow{Are external input/output transformations prescribed?}{No. The language-facing representations are specified, but transformations performed before input reaches the program or after output leaves it are outside this boundary.}
\CornellNoteRow{Are maximum program size and machine capacity fixed?}{No. The point at which a particular system cannot accommodate a program or its data is not standardized. Separate minimum translation limits constrain conforming implementations.}
\CornellNoteRow{Is a complete host-system specification provided?}{No. The clause does not attempt to define every minimum property of a data-processing system capable of supporting a conforming implementation.}
\end{CornellNotesTable}

\begin{CornellSummaryBox}{Section Summary}
A scope exclusion is not permission to violate later requirements. It identifies a mechanism or system property for which the source does not prescribe a universal design.
\end{CornellSummaryBox}



\section{Using Scope During Technical Review}
\begin{CornellNotesTable}
\CornellNoteRow{How should a disputed claim be triaged?}{First ask whether the claim concerns language meaning, observable behavior, an implementation limit, or an external mechanism. Only the first three are likely to be resolved directly by the normative clauses.}
\CornellNoteRow{What is a common category error?}{Treating a familiar compiler or operating-system convention as a language guarantee. Widely observed behavior remains nonportable unless a normative rule or documented implementation choice supports it.}
\CornellNoteRow{How does scope connect to later clauses?}{Conformance determines who must satisfy requirements; environment defines translation and execution models; language and library clauses supply the detailed rules.}
\end{CornellNotesTable}

\begin{CornellSummaryBox}{Section Summary}
Scope is the first filter for language-lawyer reasoning: determine whether the question belongs to the portable language contract before searching for a detailed rule.
\end{CornellSummaryBox}



\section{Programmer and Implementation Checklist}
\begin{itemize}
  \item[$\square$] Classify the claim as language syntax, constraint, semantics, data representation, implementation limit, or an external mechanism.
  \item[$\square$] Do not infer a portable guarantee from one compiler, ABI, operating system, or processor.
  \item[$\square$] Separate standardized input/output representation from transformations performed outside the program.
  \item[$\square$] Do not confuse minimum implementation limits with the capacity of a specific system.
  \item[$\square$] When the issue is in scope, follow the dependency chain into the conformance, environment, language, or library clause.
  \item[$\square$] When the issue is out of scope, identify the implementation or host documentation needed to resolve it.
\end{itemize}

\begin{CornellSummaryBox}{Clause Synthesis}
Clause 1 defines the perimeter of the portable programming contract. It covers program representation, language rules, program-facing data representations, and implementation restrictions. It excludes concrete toolchain mechanisms, invocation procedures, external transformations, and absolute machine capacity. A sound review therefore begins by placing each claim on the correct side of this boundary.
\end{CornellSummaryBox}



\clearpage
\section{Self-Test Questions}
\begin{enumerate}
  \item What is the principal purpose of the specification boundary?
  \item Name three categories that fall within scope.
  \item Does the source prescribe how a compiler and linker must be internally implemented?
  \item Why is machine capacity different from an implementation limit?
  \item A compiler accepts a platform-specific calling convention. Does Clause 1 make that convention portable?
  \item A program depends on how a shell converts command-line text before startup. Which boundary question applies?
\end{enumerate}

\clearpage
\section{Answer Key}
\begin{enumerate}
  \item It establishes a common interpretation of programs so that programmers and implementers can reason about portability across different systems.

  \item Examples include program representation, syntax and constraints, semantic rules, input/output data representation, and implementation restrictions or limits.

  \item No. It constrains translation results and behavior but does not prescribe the mechanism used to transform source into an executable form.

  \item A machine-capacity ceiling is specific to a system. An implementation limit is part of the documented restrictions and must still satisfy applicable minimum requirements.

  \item No. Acceptance by one implementation does not move a platform-specific mechanism into the portable language contract.

  \item The external transformation and invocation mechanisms are outside scope unless another clause imposes a relevant program-facing requirement. Host documentation is therefore needed.
\end{enumerate}

\section{Key-Term Recap}
\begin{CornellNotesTable}
\toprule
Term & Working meaning \\
\midrule
\endfirsthead
\toprule
Term & Working meaning (continued) \\
\midrule
\endhead
Scope & The defined boundary of subjects governed by the source material. \\
Program representation & The standardized form in which program text and language elements are expressed. \\
Semantic rule & A rule governing the interpretation or behavior of a valid construct. \\
Implementation limit & A restriction imposed by an implementation, subject to required minima and documentation rules. \\
External mechanism & A translation, invocation, or data-transformation process not itself prescribed by the language contract. \\
Portability & The ability to preserve intended program meaning across conforming implementations without relying on unsupported assumptions. \\
\bottomrule
\end{CornellNotesTable}

\begin{CornellSummaryBox}{One-Sentence Takeaway}
Resolve portability questions by first separating the standardized meaning of a program from the toolchain and host mechanisms that the scope intentionally leaves open.
\end{CornellSummaryBox}
\end{document}
